% ============================================================================
\section{GRASP}\label{sec:grasp}
% ============================================================================

\subsection{Idea}

GRASP (\emph{Greedy Randomized Adaptive Search Procedure}, Feo y Resende, 1989) nació
precisamente para un problema de cubrimiento difícil. Su idea es la más sencilla de las
cinco: en lugar de intentar escapar de un óptimo local, \textbf{se empieza de nuevo} desde
otra solución inicial, muchas veces, y se guarda la mejor. Cada iteración tiene dos fases:
\begin{enumerate}
  \item \textbf{Fase constructiva}: se construye una solución partiendo de cero.
  \item \textbf{Fase de búsqueda local}: se mejora esa solución.
\end{enumerate}
Para que cada iteración dé una solución distinta, la construcción no puede ser el voraz
puro (que da siempre la misma). Las tres palabras del nombre describen cómo se construye
(apuntes):
\begin{itemize}
  \item \textbf{Greedy}: en cada paso se incorpora un elemento que parece adecuado según una
        medida voraz;
  \item \textbf{Randomized}: no se elige siempre el mejor según esa medida, sino que se sortea
        entre varios ``buenos'' candidatos; por ejemplo, entre los $k$ mejores;
  \item \textbf{Adaptive}: los candidatos se reevalúan cada vez que se incorpora uno.
\end{itemize}
El conjunto de buenos candidatos entre los que se sortea se llama \textbf{lista restringida
de candidatos} (LRC; en inglés RCL). Una descripción completa y de acceso libre del método y
sus variantes es~\cite{resende}.

\subsection{Esquema de los apuntes}

\begin{apuntes}[Esquema de los apuntes: GRASP, algoritmo general]
\begin{tabbing}
\quad\=\kill
Repetir \{\\
\> Construir una solución $x\in S$ de forma greedy, aleatorizada y adaptativa\\
\> Aplicar una búsqueda local a la solución $x$\\
\> Actualizar la mejor solución conocida $x^\ast$\\
\> Actualizar la información relativa a la memoria\\
\} Mientras no se haya alcanzado el criterio de parada
\end{tabbing}
\end{apuntes}

La cuarta línea (``memoria'') distingue este esquema del GRASP básico, en el que cada
iteración es independiente de las anteriores. Los apuntes dedican una transparencia a esa
\emph{memoria y guía inteligente de la construcción}: se toma como referencia un
\textbf{conjunto élite} con un número reducido de soluciones buenas y diversas, y el criterio
de construcción pasa a ser
\[
   \lambda\,[\text{greedy}] \;+\; (1-\lambda)\,[\text{memoria}] ,
\]
con $\lambda$ grande en las primeras iteraciones (aún no se sabe nada: manda el criterio
voraz), intermedio y bajando después, y pequeño en las últimas (manda la memoria).

\subsection{Adaptación al problema de cubrimiento}

\paragraph{Fase constructiva.} Es la función \cod{completar} de la
sección~\ref{sec:completar}, partiendo de la solución vacía, con el criterio del cociente
$g_j/c_j$ y una LRC de tamaño \cod{lrc}. Es \emph{greedy} (cociente), \emph{randomized}
(sorteo dentro de la LRC) y \emph{adaptive} (las ganancias $g_j$ se actualizan tras cada
inserción).

\begin{decision}[Decisión: LRC por cardinalidad]
Hay dos formas de definir la LRC: por \emph{cardinalidad} (las $k$ mejores candidatas) o por
\emph{valor} (todas las que estén a menos de un cierto porcentaje de la mejor). Se usa la
primera, que es el ejemplo de los apuntes (``se elige aleatoriamente entre los $k$ mejores
elementos''): el parámetro es un entero con un significado inmediato y la cantidad de azar es
la misma en todos los pasos de la construcción.
\end{decision}

\paragraph{Memoria.} Se mantiene un conjunto élite con las \cod{elite} mejores soluciones
\emph{distintas} encontradas. A partir de él se calcula, para cada columna,
\[
   \text{memoria}[j] \;=\; \text{fracción de soluciones élite que contienen la columna } j
   \ \in[0,1].
\]
Una columna que aparece en todas las soluciones buenas tiene memoria 1; una que no aparece en
ninguna, 0. La puntuación con la que se forma la LRC es
\begin{equation}\label{eq:grasp}
   \text{punt}(j) \;=\; \lambda\cdot\frac{g_j/c_j}{\max_k\, g_k/c_k}
                 \;+\;(1-\lambda)\cdot\text{memoria}[j].
\end{equation}

\begin{decision}[Decisión: normalizar el término voraz]
El cociente $g_j/c_j$ puede valer 0.3 o 30 según el momento de la construcción, mientras que
la memoria está siempre entre 0 y 1. Si se sumaran tal cual, $\lambda$ no significaría nada.
Dividiendo por el mejor cociente del momento, el término voraz también queda entre 0 y 1 y
$\lambda$ es de verdad ``el peso del criterio voraz''.
\end{decision}

\begin{decision}[Decisión: $\lambda$ decrece linealmente con el tiempo]
Siguiendo las tres fases que describen los apuntes, $\lambda$ empieza en $\lambda_0=1$ y baja
en línea recta hasta $\lambda_f$ cuando se agota el tiempo:
$\lambda(t)=\lambda_0+(\lambda_f-\lambda_0)\,t/t_{\max}$. Se fija $\lambda_0=1$ porque en la
primera iteración el conjunto élite está vacío y solo puede guiar el criterio voraz. El
único valor que queda por calibrar es $\lambda_f$; con $\lambda_f=1$ la memoria no se usa
nunca y el algoritmo es el GRASP básico, lo que permite medir si la memoria aporta algo.
\end{decision}

\paragraph{Actualización del conjunto élite.} Tras cada iteración, la solución obtenida
entra en el élite si no está ya (no se admiten repetidas: el élite debe ser diverso) y,
cuando está lleno, solo si mejora a la peor, a la que sustituye. Es el compromiso ``calidad
frente a diversidad'' que mencionan los apuntes.

\paragraph{Fase de búsqueda local.} El VND de la sección~\ref{sec:vnd} con \cod{grent}
entornos. Los apuntes advierten que hay que decidir ``dónde es más interesante realizar un
mayor gasto computacional: construcción o mejora''. Con dos entornos las iteraciones son más
baratas y se hacen más; con tres, cada solución queda mejor. La calibración decide.

\subsection{Parámetros}

\begin{center}
\begin{tabular}{@{}lp{11.3cm}@{}}
\toprule
Nombre & Significado y efecto \\ \midrule
\cod{lrc} & Tamaño de la LRC. Con 1 todas las construcciones serían iguales. Pequeño:
            soluciones buenas pero parecidas. Grande: soluciones variadas pero peores. \\
\cod{lamf} & Valor final $\lambda_f$ de $\lambda$. Con 1 no hay memoria (GRASP básico). Cuanto
            menor, más pesa la memoria al final. \\
\cod{elite} & Tamaño del conjunto élite. \\
\cod{grent} & Número de entornos de la búsqueda local (2 o 3). \\
\bottomrule
\end{tabular}
\end{center}

\subsection{Recorrido por el código}

Las dos fases y el cálculo de $\lambda$:

\fragmento{grasp.c}{59}{73}

Dentro de \cod{completar}, la fórmula~\eqref{eq:grasp} es una sola instrucción (la variable
\cod{vmax} es el mejor cociente del momento):

\fragmento{base.c}{527}{531}

Y la actualización de la memoria, que solo se recalcula cuando el élite cambia:

\fragmento{grasp.c}{93}{112}
